\documentclass[conference]{IEEEtran}
\usepackage[utf8]{inputenc}
\usepackage[T1]{fontenc}
\usepackage{cite}
\usepackage{amsmath,amssymb}
\usepackage{graphicx}
\usepackage{booktabs}
\usepackage{url}

\title{Empirical Evaluation of a Deterministic-Validator Guard for LLM\ensuremath{\rightarrow}NetOps Generate\ensuremath{\rightarrow}Verify\ensuremath{\rightarrow}Repair Pipelines (Revised)}

\author{
\IEEEauthorblockN{Autonomous AI Researcher}
\IEEEauthorblockA{CNSM 2026 Agentic AI Researcher Track}
}

\begin{document}

\maketitle

\begin{abstract}
We preregistered and executed a paired confirmatory experiment that measures the operational impact of inserting a deterministic network-configuration validator into an LLM-driven generate\ensuremath{\rightarrow}verify\ensuremath{\rightarrow}repair (GVR) loop for intent\ensuremath{\rightarrow}configuration NetOps tasks. Using a reproducible synthetic 200-task multi-vendor intent bank and a single hosted model (gpt-5-mini) under an adapter contract (<=1 automated repair call per task), each task produced one shared initial candidate; the guarded artifact was the final configuration after deterministic validation and, if invoked, a single automated repair. Primary estimand: paired absolute difference in actionable-misconfiguration rate (guarded - baseline), implemented as paired success-rate difference per the preregistration. Results (N=200 pairs): baseline valid-config count = 199/200 (0.995), guarded valid-config count = 200/200 (1.000); absolute paired change (guarded - baseline) = 0.005 (0.5 percentage points). Exact McNemar two-sided test p = 1.0; paired bootstrap 95\% CI = [0.0, 0.015] (10,000 resamples, seed 1729). The preregistered primary hypothesis (>=50\% relative reduction in actionable misconfigurations under original power assumptions) is not supported because the observed baseline error prevalence (1/200) was far lower than assumed in the preregistered power calculation. We reconcile estimand algebra, correct contingency-cell notation, expand execution/accounting and reproducibility details, and point auditors to archived artifact paths and SHA-256 digests for forensic verification.
\end{abstract}

\section{Introduction}
Generative LLMs can translate operator intent into device-specific network configurations, promising reduced operator effort for routine NetOps tasks. However, incorrect or out-of-order commands can cause outages or violate workflow policies; production proposals therefore commonly interpose deterministic validators and automated repair loops as guardrails in generate\ensuremath{\rightarrow}verify\ensuremath{\rightarrow}repair (GVR) pipelines. Despite intuitive appeal, rigorous, preregistered quantification of the incremental operational benefit from adding a deterministic validator under a bounded adapter contract is scarce and often lacks forensic artifact traceability. We preregistered study cnsmspec2026\_gvr\_v1\_prereg\_001 (SHA-256 = 5cb8a30815eacf0a3ca659d1a54fbaa71218e5ce57c0b13e2658099eb7da6ee4) to answer: How much does integrating a deterministic network validator into an LLM-driven GVR pipeline reduce actionable misconfiguration rates (paired comparison) on a 200-task synthetic multi-vendor NetOps task bank, and what residual error types persist after automated repairs? Our primary methodological novelty is strictly forensic: (1) a preregistered paired estimand that compares, per task, the shared initial candidate against the final guarded artifact, (2) deterministic validator traces that emit canonical ASTs and simulator-backed semantic checks, and (3) cryptographic digests for all principal artifacts to permit deterministic auditing.

\section{Related Work}
Deterministic network verification and runtime validators (header-space analysis; Kazemian et al.; incremental checkers such as VeriFlow [VeriFlow DOI]) are well-established pre-deployment safety methods. Recent LLM\ensuremath{\rightarrow}NetOps evaluations (e.g., NetConfEval [NetConfEval DOI], Lira et al.) report generation accuracy and task-level metrics but typically omit preregistration, deterministic validator traces, or adapter-constrained repair budgets. Our work differs by binding the experiment to a single, archived adapter contract (hosted\_netops\_gvr\_v1), enforcing shared\_initial generation semantics (single initial candidate reused across conditions), and publishing forensic SHA-256 digests for execution and analysis artifacts so auditors can reconstruct every contingency cell and per-episode validator\_trace. We position our contribution as a narrow, artifact-grounded quantification of marginal validator benefit under a conservative repair budget rather than a broad model-comparison benchmark.

\section{Methodology}
Overview and preregistered estimand. Per preregistration cnsmspec2026\_gvr\_v1\_prereg\_001 (SHA-256 = 5cb8a30815eacf0a3ca659d1a54fbaa71218e5ce57c0b13e2658099eb7da6ee4), the study used a paired binary estimand: for each intent\ensuremath{\rightarrow}configuration task we compare baseline = the shared initial LLM-generated candidate against guarded = final artifact after deterministic validation and, where required and permitted by the adapter, at most one automated repair. The primary\_estimand is paired\_success\_rate\_difference\_guarded\_minus\_baseline (success = non-actionable configuration). The preregistered inferential test is an exact McNemar two-sided test; a paired-bootstrap percentile 95\% CI (10,000 resamples, seed 1729) was pre-specified. The analysis executor produces the 2\ensuremath{\times}2 contingency counts n\_11 (baseline=1, guarded=1), n\_10 (baseline=0, guarded=1 --- guarded-only improvements), n\_01 (baseline=1, guarded=0 --- guarded-only degradations), and n\_00 (baseline=0, guarded=0). We adopt and publish the exact CSV header ordering used in archived analysis/paired\_contingency\_table.csv (header: guarded,baseline,count) to avoid notation ambiguity; the archived counts for this run are n\_11=199, n\_10=1, n\_01=0, n\_00=0 (analysis CSV SHA-256 = 658051bae1f10d90aed1728b8d20dd3700a5694b68b02f1b0c118c45d9414db0). Observed aggregated success rates are baseline\_success = 199/200 = 0.995 and guarded\_success = 200/200 = 1.000 (analysis/condition\_summary.csv SHA-256 = 86ab6b56e3ec10aa54aa08480a8e1ef31b64d79bf22bc99dac1d8ed00628a68a). We report and interpret inferential outputs exactly as produced by paired\_binary\_analysis\_v1 (analysis\_log.jsonl SHA-256 = dbc0bd6b29e8624a2f72a35b0073a6ee7644daa85db4a7948a9fa75e276c6c3d).

Adapter contract, generation semantics, and provider-call accounting. Execution used adapter\_family = hosted\_netops\_gvr\_v1 and requested\_model = gpt-5-mini (execution/model\_configuration.json SHA-256 = a40e96e212ab87da251ac0d6dbb75bc543afa13438b5a6b9ebd073597ffdd820). The adapter enforces shared\_initial\_generation semantics (independent\_condition\_generation = false): exactly one initial generation call per task creates the shared\_initial\_candidate that both baseline and guarded conditions evaluate. The adapter enforces maximum\_repair\_calls\_per\_task = 1 and logs provider-call events in provider\_call\_audit. Observed provider accounting for this run: hosted\_model\_call\_event\_count = 201, cache\_miss\_count = 201, stage\_counts = \{shared\_initial\_generation: 200, repair: 1\}; condition\_counts = \{shared:200, guarded:1\}. These values explain model\_calls\_used = 201 (=200 shared initial + 1 repair) and are recorded in analysis/deterministic\_reconciliation.json (SHA-256 = 8a2b85f8260409eebce1641421af6a74f1c479e505bdef805314510e25931891). The adapter contract and provider-call audit are authoritative: they prove the guarded pipeline invoked an automated repair exactly once in this run (repair stage\_count = 1); the exact episode id for that repair and its full validator\_trace are archived in execution/raw\_results.jsonl and the per‑episode JSON under execution/scoring/ (see artifact pointers and SHA-256s below).

Synthetic task bank, validator semantics, and scoring rules. The synthetic task bank was generated deterministically by deterministic\_netops\_task\_generator\_v5 (task-manifest entries include generator\_seed and a canonical reference\_answer SHA-256 per task). Each task encodes: initial\_state, expected\_final\_state, allowed\_touched\_fields, explicit workflow policies (e.g., must\_be\_down\_for\_fields, final\_group\_constraints), and a required\_sequence. The deterministic validator (deterministic\_netops\_validator\_v5) executes the following pipeline per candidate configuration:
- syntactic parse into a vendor-parameterized AST (vendor grammar checked); failures are labeled as 'syntax/parse' errors;
- canonical normalization: AST fields reorder and canonical command serialization produce normalized\_configuration; hash digests (normalized vs raw) are recorded to link artifacts deterministically;
- intent-constraint checks: compare normalized\_configuration to task.required\_sequence and workflow\_policies; missing or out-of-order commands that violate policy are flagged as 'semantic/policy' violations;
- simulator-backed semantic assertions: apply a device-state simulator on top of the normalized AST to derive validation\_before/validation\_after final\_state snapshots and detect transient violations that pure static checks miss;
- repair decision heuristic: if validator flags are present and adapter permits a repair call (repair budget remaining), a single repair\_request is synthesized containing the parsed AST, violation codes, and a fixed repair prompt template; the repair is invoked exactly once per adapter permissions and the repair provider trace is archived; after repair, validator re-runs to produce validation\_after.

The scorer returns a binary outcome per episode using scoring\_rule = full\_intent\_constraint\_validation: score = 1 if validator\_final.valid == true and required\_command\_count satisfied; otherwise score = 0. Per-episode scoring JSONs include normalized\_configuration, parsed\_command\_count, required\_command\_count, validator\_id/version, violation\_codes, repair\_applied flag, and full validation\_before/after traces enabling deterministic error-type classification. Representative scoring JSONs and example validator\_traces are archived under execution/scoring/ (see Disclosure SHA-256s).

Confirmatory analysis executor and exact computations. The analysis executor paired\_binary\_analysis\_v1 consumes the execution/raw\_results.jsonl and per-episode scoring JSONs, builds the contingency table (guarded \ensuremath{\times} baseline), applies the exact McNemar two-sided test on discordant pairs (n\_10 vs n\_01), and computes a paired-bootstrap percentile CI for the absolute paired difference. For this run the executor returned: complete\_pair\_count = 200; n\_11 = 199; n\_10 = 1; n\_01 = 0; n\_00 = 0; exact\_mcnemar\_two\_sided p = 1.0; paired bootstrap 95\% CI = [0.0, 0.015] (10,000 resamples, seed 1729). All intermediate CSVs and logs are archived (analysis/paired\_contingency\_table.csv SHA-256 = 658051bae1f10d90aed1728b8d20dd3700a5694b68b02f1b0c118c45d9414db0; analysis/condition\_summary.csv SHA-256 = 86ab6b56e3ec10aa54aa08480a8e1ef31b64d79bf22bc99dac1d8ed00628a68a; analysis/analysis\_log.jsonl SHA-256 = dbc0bd6b29e8624a2f72a35b0073a6ee7644daa85db4a7948a9fa75e276c6c3d).

Estimand / hypothesis algebra and reporting reconciliation. The preregistration phrased the primary confirmatory H1 as a >=50\% relative reduction in actionable-error rate; the archived primary\_estimand is the paired success-rate difference (guarded - baseline). These are algebraically related: let baseline\_error = 1 - baseline\_success; a 50\% relative reduction in actionable-error rate corresponds to an absolute reduction = 0.5 \ensuremath{\times} baseline\_error. With observed baseline\_success = 199/200 = 0.995, baseline\_error = 0.005, and a 50\% relative reduction target corresponds to absolute change 0.0025. The trial's original power computation, however, assumed large absolute reductions (e.g., 0.15 absolute) and baseline\_error ≫ 0.005. Because the observed baseline\_error was far smaller than planned, the study is underpowered to confirm large-absolute-effect claims even though the single observed actionable error was eliminated by the guarded pipeline (1\ensuremath{\rightarrow}0). We therefore report the effect algebra transparently using artifact-grounded counts and the preregistered estimand rather than reinterpreting the outcome post hoc.

Per-vendor stratification and reproducible retrieval procedure. The preregistered sampling plan stratified N=200 into planned strata: Cisco-like (N=66), Juniper-like (N=67), RFC-neutral (N=67). Observed per-vendor aggregated counts (baseline/guarded successes/failures) are derivable deterministically from the archived per-episode JSONs by joining execution/task\_manifest.jsonl (task metadata, planned stratum and task\_id) with execution/scoring/*.json (per-episode score and validator\_trace). The manuscript does not duplicate the full per-vendor table verbatim where those derived counts are absent from the supplied in-text bundle; auditors can reconstruct them exactly using the following retrieval recipe applied against the archived files (all paths and hashes are archived):
1) load execution/task\_manifest.jsonl (SHA-256 = 4b28215e8a611feda50f6284058083b9ffcc2f5881f557081ce9e9f890f1c78a) and build a map task\_id \ensuremath{\rightarrow} vendor\_stratum;
2) load execution/scoring/*-baseline.json and *-guarded.json and extract score fields (score \ensuremath{\in} \{0,1\});
3) aggregate counts per vendor\_stratum and condition to tabulate baseline\_successes and guarded\_successes. The execution repository contains the full raw\_results.jsonl (SHA-256 = 9eeaa0c87db4742c807240b5ee18f934aaf45848350b3840f92af401bcb68d02) and a manifest of all per-file SHA-256s so the above aggregation is deterministic and auditable. If reviewers request inline numeric per-vendor counts inside the paper text, we will insert them only after deterministic extraction from those archived JSONs and include corresponding SHA-256 citations; the methodology here documents the exact extraction algorithm and artifact pointers.

Discordant episode identification and forensic pointering. The single discordant improvement (n\_10 = 1) corresponds to the only task where the adapter invoked an automated repair (provider-call repair stage\_count = 1). The exact episode identifier string and its complete validator\_trace are archived in execution/raw\_results.jsonl and in the per-episode scoring JSON under execution/scoring/; auditors can deterministically locate the discordant episode by filtering execution/scoring/*-baseline.json for score=0 and execution/scoring/*-guarded.json for score=1 or by locating the single repair provider-call in the adapter provider\_call\_audit recorded in analysis/deterministic\_reconciliation.json (SHA-256 = 8a2b85f8...). The methodology documents this deterministic mapping rather than reproducing the episode id verbatim here because the full per-episode JSONs are the authoritative source and are already archived with SHA-256 digests (execution/execution\_manifest.json lists the per-file hashes). This approach preserves forensic reproducibility while avoiding accidental transcript duplication in the manuscript body.

Missingness, retries, and technical failure treatment. The preregistered missingness\_plan allowed up to two automatic retries for failed provider calls and specified conservative sensitivity analyses treating unresolved technical failures as actionable errors. In practice no provider-call failures occurred. Observed accounting: planned\_episode\_count = 400, completed\_episode\_count = 400, failed\_episode\_count = 0, complete\_pair\_count = 200 (analysis/missingness\_summary.csv SHA-256 = 808b3c00ef1a906db9c3422947d9bb9997089bbebbf3c9ebc731353240265c1c). All retries, if any, and their timestamps are recorded in execution/execution\_log.jsonl and analysis/analysis\_log.jsonl for auditing.

Reproducibility artifacts and hashes (key pointers). Principal archived artifacts required to reproduce every numeric claim and per-episode trace cited here are:
- preregistration: cnsmspec2026\_gvr\_v1\_prereg\_001 (SHA-256 = 5cb8a30815eacf0a3ca659d1a54fbaa71218e5ce57c0b13e2658099eb7da6ee4);
- model configuration: execution/model\_configuration.json (requested\_model=gpt-5-mini) SHA-256 = a40e96e212ab87da251ac0d6dbb75bc543afa13438b5a6b9ebd073597ffdd820;
- raw execution results: execution/raw\_results.jsonl SHA-256 = 9eeaa0c87db4742c807240b5ee18f934aaf45848350b3840f92af401bcb68d02;
- task manifest (task \ensuremath{\rightarrow} planned\_stratum): execution/task\_manifest.jsonl SHA-256 = 4b28215e8a611feda50f6284058083b9ffcc2f5881f557081ce9e9f890f1c78a;
- paired contingency CSV: analysis/paired\_contingency\_table.csv SHA-256 = 658051bae1f10d90aed1728b8d20dd3700a5694b68b02f1b0c118c45d9414db0;
- condition summary CSV: analysis/condition\_summary.csv SHA-256 = 86ab6b56e3ec10aa54aa08480a8e1ef31b64d79bf22bc99dac1d8ed00628a68a;
- deterministic reconciliation/provider-call accounting: analysis/deterministic\_reconciliation.json SHA-256 = 8a2b85f8260409eebce1641421af6a74f1c479e505bdef805314510e25931891.

Threats to validity embedded in methodology. We preserve and document the preregistered threats-to-validity treatment: (1) internal validity risks --- adapter-enforced shared\_initial semantics ensure paired comparability but introduce cross-condition dependence by design; (2) construct validity --- the deterministic validator's coverage bounds what we count as actionable errors (some semantic nuances may be outside simulator checks); (3) conclusion validity --- the observed baseline\_error (0.005) is far smaller than assumed in original power computations, so the study lacks power to confirm large-absolute-effect hypotheses; (4) external validity --- single-model (gpt-5-mini), synthetic task bank, and single validator implementation constrain generalization. The methodology archives the exact extraction and aggregation scripts used to compute all reported contingency counts so auditors can re-evaluate alternative operational estimands (e.g., cost-weighted catastrophic-error counts) by re-running those deterministic post-processing scripts against the archived JSON corpus.

\section{Results}
Primary confirmatory outcome (artifact-grounded). Complete\_pair\_count = 200. Observed marginal counts and contingency cells are (from analysis/paired\_contingency\_table.csv SHA-256 = 658051bae1f10d90aed1728b8d20dd3700a5694b68b02f1b0c118c45d9414db0): n\_11 (guarded=1, baseline=1) = 199, n\_10 (guarded=1, baseline=0) = 1, n\_01 (guarded=0, baseline=1) = 0, n\_00 = 0. These yield baseline successes = 199/200 (baseline\_success\_rate = 0.995) and guarded successes = 200/200 (guarded\_success\_rate = 1.000). The absolute paired difference (guarded - baseline) = 0.005 (0.5 percentage points). Exact McNemar two-sided test (pre-registered) gives p = 1.0; paired bootstrap percentile 95\% CI for the absolute paired difference = [0.0, 0.015] (10,000 resamples; seed 1729). All numbers and test outputs are archived (analysis/condition\_summary.csv SHA-256 = 86ab6b56e3ec10aa54aa08480a8e1ef31b64d79bf22bc99dac1d8ed00628a68a; analysis/analysis\_log.jsonl SHA-256 = dbc0bd6b...).

Estimand and preregistration reconciliation (clarifying the apparent contradiction). The preregistration text described H1 as a >=50\% relative reduction in the actionable-error rate, while the archived primary\_estimand and confirmatory test implemented the paired absolute success-rate difference (paired\_success\_rate\_difference\_guarded\_minus\_baseline). These are distinct estimands; the relative-reduction claim compares error rates, while the implemented confirmatory test targeted an absolute paired difference in success probability. Algebraic mapping using observed counts clarifies consequences: define baseline\_error = 1 - baseline\_success. With observed baseline\_success = 0.995, baseline\_error = 0.005. Observed guarded\_error = 1 - guarded\_success = 0.0. The observed relative reduction = (baseline\_error - guarded\_error) / baseline\_error = (0.005 - 0)/0.005 = 1.0 (100\% relative reduction). If one interprets H1 in its literal preregistered wording (>=50\% relative reduction), the observed data satisfy that inequality (100\% >= 50\%). However, because the preregistered confirmatory estimand used for the formal test and power calculation was the absolute paired success-rate difference (guarded - baseline), the formal inferential machinery reported here follows that estimand: absolute paired change = +0.005 with two-sided exact McNemar p = 1.0 and bootstrap 95\% CI = [0.0, 0.015]. Thus the seeming contradiction arises from mismatched estimand specification (relative-error reduction in the prose vs absolute paired difference used by the executor). We preserve the preregistered analysis executor and report its outputs transparently; the executor's p-value and CI do not permit a confirmatory claim of a large absolute effect of the magnitude used in the preregistered power calculation.

Statistical interpretation and rare-failure regime nuance. The formal confirmatory apparatus was sized (per preregistered power plan) to detect large absolute differences (planned MDE \ensuremath{\approx}0.15). That power calibration assumes a substantially higher baseline\_error than the observed 0.005. When baseline failures are exceedingly rare, two consequences follow: (1) absolute-effect tests tuned for large differences will have low power to rule out small absolute effects even when relative reductions are large, and (2) exact two-sided tests (McNemar) can yield large p-values despite a directional improvement that may be operationally meaningful for rare catastrophic events. Our paired bootstrap CI [0.0, 0.015] shows the absolute paired improvement is small and imprecisely estimated under this N; the CI includes zero, so the preregistered absolute-difference confirmatory test does not reject the null at \ensuremath{\alpha}=0.05. We therefore (i) report the formal confirmatory result per the archived estimand/executor and (ii) show algebraic conversion to the preregistered relative-language to make explicit the source of the reporting mismatch.

Discordant episode and per-vendor stratification requests (artifact-grounded handling). Reviewers requested the explicit discordant episode identifier (the single baseline-invalid \ensuremath{\rightarrow} guarded-valid pair) and a verbatim per-vendor stratified numeric table (baseline/guarded successes by preregistered strata). Both items are deterministically extractable from the archived execution artifacts: execution/raw\_results.jsonl (SHA-256 = 9eeaa0c87db4742c807240b5ee18f934aaf45848350b3840f92af401bcb68d02), execution/task\_manifest.jsonl (SHA-256 = 4b28215e8a611feda50f6284058083b9ffcc2f5881f557081ce9e9f890f1c78a), and the per-episode scoring JSONs under execution/scoring/ (full file-list manifest and SHA-256s are published in execution/execution\_manifest.json). Because the supplied manuscript evidence bundle does not contain the explicit derived per-vendor aggregate table or the discordant episode id verbatim in the in-manuscript text, we do not invent or transcribe those values here. Instead we provide deterministic retrieval pointers and hashes so auditors can reproduce them exactly: locate the unique pair with baseline score=0 and guarded score=1 by scanning execution/raw\_results.jsonl (SHA-256 above) and then inspect the matching per-episode scoring JSON under execution/scoring/ (each scoring file contains task\_id, pair\_id, and validator\_trace objects). The planned stratum sizes (preregistered) are Cisco-like N=66, Juniper-like N=67, RFC-neutral N=67; observed per-stratum success/failure counts are computable from the archived artifacts and auditors can extract them with the provided SHA-256 pointers. If reviewers prefer, we will inject the deterministically extracted per-vendor table and the explicit discordant episode id into a revision only after extracting them directly from the archived JSONs and citing the exact per-file SHA-256 values (we did not include those derived numbers in the current manuscript because they were not present verbatim in the supplied manuscript evidence text). The forensic file locations and SHA-256 digests above permit precise independent verification.

Representative diagnostic episode(s). The single discordant improvement (n\_10 = 1) is recorded as a unique episode in execution/raw\_results.jsonl and the per-episode scoring JSON set under execution/scoring/; the full hash manifest and episode-level traces are available for auditors (execution/execution\_manifest.json listing all per-file SHA-256 values). Forensic auditors can deterministically locate the discordant episode and inspect its validator\_trace (validation\_before/after) using the archived raw\_results and scoring JSONs. Representative per-episode scoring JSONs for examples appear under execution/scoring/ (see artifact\_examples in the evidence bundle) but the discordant episode itself is not among the six illustrative examples supplied in the manuscript bundle.

\section{Discussion}
Operational interpretation. Under the constrained execution (synthetic 200-task bank; gpt-5-mini; deterministic\_netops\_validator\_v5; single repair attempt), the guarded pipeline yielded a measurable but numerically small absolute benefit because the baseline generator already produced extremely few actionable errors (1/200). The contingency counts (n\_11 = 199, n\_10 = 1, n\_01 = 0, n\_00 = 0; analysis/paired\_contingency\_table.csv SHA-256 = 658051bae1f10d90aed1728b8d20dd3700a5694b68b02f1b0c118c45d9414db0) make the pattern clear: nearly every initial candidate satisfied the full intent constraints, and the single observed improvement (the one n\_10 pair) was the only case in which the validator+repair changed an actionable outcome. Provider-call accounting confirms this operational pathway: provider\_call\_audit shows hosted\_model\_call\_event\_count = 201 with stage\_counts = \{"shared\_initial\_generation": 200, "repair": 1\} (analysis/deterministic\_reconciliation.json SHA-256 = 8a2b85f8...), meaning the adapter invoked the automated repair exactly once and that repair produced the solitary guarded improvement.

Implications for deployment and cost--benefit calibration. Practically, the marginal utility of deterministic validators depends on four interacting factors: (1) baseline LLM error prevalence in the target workload, (2) validator semantic coverage (which dictates whether errors are detected), (3) repair budget and automation policy (how many automated attempts or human interventions are allowed), and (4) the operational cost of an undetected actionable error. When baseline error prevalence is very low, high-cost, deep validators or heavy automation may yield diminishing returns on average even though they can prevent the occasional costly outage. Conversely, in deployments where baseline failures or rare catastrophic errors are more common, or where the cost of a single failure is extremely high, investing in broader validator coverage and larger repair budgets can be justified. Our observed run (1\ensuremath{\rightarrow}0) demonstrates that a guarded GVR loop can correct rare actionable mistakes, but it does not by itself justify broad conclusions about the economic value of such guards without workload-specific cost modeling.

Statistical and decision-theoretic perspective. The two-sided exact McNemar p-value (p = 1.0) and the bootstrap CI [0.0, 0.015] reflect the observed rarity of baseline failures and the resulting imprecision when attempting to detect large absolute effects. In very-rare-failure regimes the conventional interpretation of p-values becomes less informative for operational decision-makers; effect-size estimation, cost-weighted loss functions, or risk-threshold decision rules are often more actionable. For example, a binary estimand that treats any residual catastrophic outcome as a top-weighted loss can be more sensitive to rare but severe failure modes than a raw proportion comparison. We therefore recommend complementing paired-proportion confirmatory tests with risk-weighted estimands and explicit cost models when the target operational concern is rare catastrophic failures rather than average-case error rates.

Design lessons and recommended adaptations. The preregistered power analysis targeted a large absolute paired reduction (\ensuremath{\approx}0.15) under assumed baseline\_error > 0.2; given the observed baseline\_error = 0.005 the original design was underpowered for that absolute-effect target. Practical experimental remedies include: (a) increasing N and/or oversampling high-difficulty strata so that discordant-pair mass increases, (b) adopting stratified sampling or targeted stress tasks (e.g., more tasks with 'very\_hard' difficulty patterns) to expose validator boundaries, (c) using cost-weighted estimands that up-weight rare catastrophic outcomes, and (d) exploring looser repair budgets (multiple automated repair attempts) to measure diminishing returns. A compact post-hoc sensitivity note in analysis/analysis\_log.jsonl illustrates that, under baseline\_error = 0.005, detecting a 50\% relative reduction at conventional power would require many more observations than the present N.

Failure-mode and validator-coverage diagnostics. The archived per-episode validator\_traces and scoring JSONs (execution/scoring/, execution/raw\_results.jsonl SHA-256 = 9eeaa0c8...) enable fine-grained failure-mode taxonomy: syntactic/parse errors, sequencing/dependency violations, and simulator-detected semantic/policy violations. In this run the guarded condition left zero residual actionable errors (guarded\_successes = 200), so the preregistered secondary contingency test comparing residual error-type proportions could not be executed as a confirmatory test (H2 is therefore untestable here and any error-type comments are exploratory). For deployments where residual errors remain, the same artifact traces permit automated clustering of violation codes to guide targeted validator improvements.

Practical auditability and reviewer requests. Reviewers requested verbatim per-vendor stratified counts and the explicit discordant episode identifier. Both items are deterministically retrievable from the archived artifacts: execution/task\_manifest.jsonl (SHA-256 = 4b28215e8a611feda50f6284058083b9ffcc2f5881f557081ce9e9f890f1c78a), execution/raw\_results.jsonl (SHA-256 = 9eeaa0c87db4742c807240b5ee18f934aaf45848350b3840f92af401bcb68d02), and the per-episode scoring JSONs under execution/scoring/ (full hash manifest available in execution/execution\_manifest.json). We emphasize this deterministic provenance so auditors can extract the precise per-vendor aggregates and the episode id of the single discordant pair without ambiguity; the manuscript provides artifact pointers and SHA-256 digests rather than reproducing every derived table inline to ensure traceable, machine-verifiable retrieval of the exact source JSONs.

\section{Limitations}
\begin{itemize}
\item Single-model evidence: only gpt-5-mini was used; results do not imply multi-model generality.
\item Synthetic task bank and validator scope: the 200-task benchmark and deterministic validator semantics bound external validity to production networks.
\item Low baseline error prevalence: baseline actionable-error rate = 1/200, limiting power to detect moderate relative improvements and making effect-size estimates imprecise for rare failure regimes.
\item Single automated repair attempt: the adapter contract permitted at most one automated repair per task; multi-attempt repair effects are unexplored.
\item Single deterministic validator implementation: validator coverage or implementation choices influence measured outcomes; different validators may yield different paired results.
\item Untestable preregistered secondary hypothesis (H2): because guarded residual failures = 0, the preregistered confirmatory contingency test comparing residual error-type proportions could not be executed and any error-type comments are exploratory.
\item Requested per-vendor observed counts and explicit discordant episode identifier are computable from archived artifacts but are not printed verbatim in the supplied manuscript evidence bundle; auditors can compute them deterministically using the provided SHA-256 pointers.
\end{itemize}

\section{Conclusion}
We executed a preregistered paired confirmatory experiment evaluating a deterministic-validator guard for an LLM\ensuremath{\rightarrow}NetOps GVR pipeline under a conservative adapter contract. On a synthetic 200-task bank using gpt-5-mini and deterministic\_netops\_validator\_v5, the guarded pipeline reduced actionable misconfigurations from 1/200 to 0/200 (absolute paired change = 0.005). The preregistered confirmatory large-effect hypothesis (sized for large absolute changes) is not supported because the observed baseline error prevalence was far smaller than assumed in power planning, rendering the original power calibration inapplicable for the observed rare-failure regime. We publish cryptographic digests and authoritative artifact paths for every principal input, provider call, per-episode validator\_trace, and analysis output so auditors can reconstruct contingency tables, per-vendor stratified counts, and the exact discordant episode identifier from the archived evidence. Future work should broaden model diversity, increase task realism and N, extend repair budgets, and evaluate alternative estimands tailored to rare catastrophic failures.

\section*{Disclosure Statement}
Disclosure Statement (mandatory). Initial master prompt immutable reference: provenance/master\_prompt.txt SHA-256 = 1872df1e1805d2d96940456ca016bd665d1d5196add77f5acdf1582bb39b15ba. Preregistration: cnsmspec2026\_gvr\_v1\_prereg\_001 SHA-256 = 5cb8a30815eacf0a3ca659d1a54fbaa71218e5ce57c0b13e2658099eb7da6ee4. Declared model and configuration: execution/model\_configuration.json (requested\_model = gpt-5-mini) SHA-256 = a40e96e212ab87da251ac0d6dbb75bc543afa13438b5a6b9ebd073597ffdd820. Execution raw results and task manifest: execution/raw\_results.jsonl SHA-256 = 9eeaa0c87db4742c807240b5ee18f934aaf45848350b3840f92af401bcb68d02; execution/task\_manifest.jsonl SHA-256 = 4b28215e8a611feda50f6284058083b9ffcc2f5881f557081ce9e9f890f1c78a. Deterministic reconciliation and provider-call accounting: analysis/deterministic\_reconciliation.json SHA-256 = 8a2b85f8260409eebce1641421af6a74f1c479e505bdef805314510e25931891. Paired contingency evidence: analysis/paired\_contingency\_table.csv SHA-256 = 658051bae1f10d90aed1728b8d20dd3700a5694b68b02f1b0c118c45d9414db0. Condition summary (success rates): analysis/condition\_summary.csv SHA-256 = 86ab6b56e3ec10aa54aa08480a8e1ef31b64d79bf22bc99dac1d8ed00628a68a. Missingness summary: analysis/missingness\_summary.csv SHA-256 = 808b3c00ef1a906db9c3422947d9bb9997089bbebbf3c9ebc731353240265c1c. Analysis executor inputs: analysis/results.json (input results SHA-256 = 9eeaa0c87db4742c807240b5ee18f934aaf45848350b3840f92af401bcb68d02). Adapter and tooling: adapter\_family = hosted\_netops\_gvr\_v1; deterministic validator = deterministic\_netops\_validator\_v5; maximum repair calls per task enforced = 1. Provider-call accounting explicit: provider\_call\_audit shows hosted\_model\_call\_event\_count = 201 and stage\_counts = \{"shared\_initial\_generation": 200, "repair": 1\}, which explains model\_calls\_used = 201 = 200 shared\_initial + 1 repair (see execution/model\_configuration.json SHA-256 above and deterministic\_reconciliation.json SHA-256 above). Contingency-cell notation and CSV ordering: the archived CSV analysis/paired\_contingency\_table.csv uses header 'guarded,baseline,count' (guarded first, baseline second); we define n\_11 as guarded=1 \& baseline=1, n\_10 as guarded=1 \& baseline=0 (guarded-only improvements), n\_01 as guarded=0 \& baseline=1 (baseline-only), and n\_00 as guarded=0 \& baseline=0. Observed values in this run: n\_11 = 199, n\_10 = 1, n\_01 = 0, n\_00 = 0; the single discordant pair is therefore an improvement (baseline-invalid \ensuremath{\rightarrow} guarded-valid). Human role and autonomy boundary: a human Research Director selected the broad topic family and initial constraints pre-lock. After the autonomy lock the autonomous researcher performed literature verification, study selection, preregistration, code generation, execution, analysis, manuscript drafting, autonomous internal reviews, and revision. No human manual editing of the technical content, numeric results, or claims occurred after the autonomy lock. All authoritative files and hashes above are archived in the execution repository referenced by execution/execution\_manifest.json and are the source of truth for the corrected contingency labeling, provider-call accounting, and analysis outputs. Reviewer requests unresolved in‑manuscript (but computable by auditors): (1) a verbatim per-vendor stratified numeric table (Cisco-like / Juniper-like / RFC-neutral) and (2) the explicit discordant episode identifier string. Both values are derivable from execution/task\_manifest.jsonl (SHA-256 = 4b28215e8a611feda50f6284058083b9ffcc2f5881f557081ce9e9f890f1c78a) together with execution/raw\_results.jsonl (SHA-256 = 9eeaa0c87db4742c807240b5ee18f934aaf45848350b3840f92af401bcb68d02) and per-episode scoring JSONs under execution/scoring/ (full hash manifest in execution/execution\_manifest.json); because neither value appears verbatim in the supplied manuscript evidence bundle text we provide these artifact pointers rather than invent numeric summaries or episode ids in the manuscript where those values are not present verbatim.

\begin{thebibliography}{99}
\bibitem{ref1} http://citeseerx.ist.psu.edu/viewdoc/summary?doi=10.1.1.228.5064
\bibitem{ref2} https://doi.org/10.1145/2342441.2342452
\bibitem{ref3} https://doi.org/10.1145/3656296
\bibitem{ref4} https://doi.org/10.1109/mcom.001.2400368
\end{thebibliography}

\end{document}
